% Shared LaTeX book preamble — XeLaTeX

% 기본 패키지
\usepackage{kotex,fontspec,mathtools,unicode-math,amsthm}
\usepackage{setspace,indentfirst,enumitem}
\usepackage{graphicx,booktabs,array,longtable,diagbox,listings,makeidx}
\usepackage{tikz}
\usetikzlibrary{positioning,calc,arrows.meta,patterns}
\usepackage{xcolor,titlesec,titletoc,fancyhdr}
\usepackage[most]{tcolorbox}

% 글꼴
% 운영체제의 글꼴 등록 상태나 절대경로에 의존하지 않고,
% TeX Live와 저장소가 제공하는 실제 글꼴 파일을 직접 사용한다.
\setmainfont{STIXTwoText-Regular.otf}[
  BoldFont=STIXTwoText-Bold.otf,
  ItalicFont=STIXTwoText-Italic.otf,
  BoldItalicFont=STIXTwoText-BoldItalic.otf
]
\setmonofont{lmmonolt10-regular.otf}[
  BoldFont=lmmonolt10-bold.otf,
  ItalicFont=lmmonolt10-oblique.otf,
  BoldItalicFont=lmmonolt10-boldoblique.otf
]
\setmathfont{STIXTwoMath-Regular.otf}

% KoPubWorld Pro OTF 6개는 fonts/에 둔다.
\IfFileExists{fonts/KoPubWorld Batang_Pro Light.otf}{
  \setsansfont{KoPubWorld Dotum_Pro}[
    Path=fonts/,
    Extension=.otf,
    UprightFont={* Light},
    BoldFont={* Bold},
    ItalicFont={* Light},
    BoldItalicFont={* Bold}
  ]
  \setmainhangulfont{KoPubWorld Batang_Pro}[
    Path=fonts/,
    Extension=.otf,
    UprightFont={* Light},
    BoldFont={* Bold},
    ItalicFont={* Light},
    BoldItalicFont={* Bold},
    SmallCapsFont={* Light}
  ]
  \setsanshangulfont{KoPubWorld Dotum_Pro}[
    Path=fonts/,
    Extension=.otf,
    UprightFont={* Light},
    BoldFont={* Bold},
    ItalicFont={* Light},
    BoldItalicFont={* Bold}
  ]
  \setmonohangulfont{KoPubWorld Dotum_Pro}[
    Path=fonts/,
    Extension=.otf,
    UprightFont={* Light},
    BoldFont={* Bold},
    ItalicFont={* Light},
    BoldItalicFont={* Bold}
  ]
}{
  \PackageError{latexbook}{KoPubWorld font files not found}
    {Upload the six KoPubWorld Pro OTF files to the fonts directory.}
}

% 판형과 본문
\usepackage[
  paperwidth=182mm,paperheight=257mm,
  top=22mm,bottom=25mm,inner=22mm,outer=18mm,
  headsep=8mm,footskip=12mm
]{geometry}
\setlength{\parindent}{1em}
\setlength{\parskip}{0pt}
\emergencystretch=2em
\setlist{nosep,leftmargin=2em}

% 그림·표·색인
\graphicspath{{figures/}}
\makeindex
\renewcommand{\indexname}{찾아보기}
\renewcommand{\figurename}{그림}
\renewcommand{\tablename}{표}

% 흑백 인쇄용 회색조 팔레트. 의미 구분은 선종류·마커·라벨로도 반복한다.
\definecolor{BookPurple}{gray}{0.00}
\definecolor{BookBlue}{gray}{0.00}
\definecolor{BookGreen}{gray}{0.35}
\definecolor{BookOrange}{gray}{0.25}
\definecolor{BookRed}{gray}{0.15}
\definecolor{BookGray}{gray}{0.38}

% 장·절 제목
\titleformat{\part}[display]
  {\centering\sffamily\bfseries}{\Large 제 \thepart\ 부}{1.2em}{\Huge}
\titlespacing*{\part}{0pt}{42mm}{0pt}
\titleformat{\chapter}[display]
  {\sffamily\bfseries}{\Large 제 \thechapter\ 장}{1.1em}{\Huge}
\titlespacing*{\chapter}{0pt}{25mm}{11mm}
\titleformat{\section}
  {\sffamily\large\bfseries}{\thesection}{0.75em}{}
\titlespacing*{\section}{0pt}{3.0ex plus 0.8ex minus 0.4ex}{1.2ex}
\titleformat{\subsection}
  {\sffamily\normalsize\bfseries}{\thesubsection}{0.75em}{}
\titlespacing*{\subsection}{0pt}{2.4ex plus 0.6ex minus 0.3ex}{0.9ex}
\renewcommand{\contentsname}{목차}

% 머리말과 꼬리말
\pagestyle{fancy}
\fancyhf{}
\renewcommand{\headrulewidth}{0pt}
\renewcommand{\footrulewidth}{0pt}
\renewcommand{\chaptermark}[1]{\markboth{#1}{}}
\renewcommand{\sectionmark}[1]{\markright{#1}}
\fancyhead[LE]{\small\color{BookGray}--- \BookTitle}
\fancyhead[RO]{\small\color{BookGray}\BookSubtitle{} ---}
\fancyfoot[LE]{\thepage\hspace{4mm}\small\leftmark}
\fancyfoot[RO]{\small\rightmark\hspace{4mm}\thepage}
\fancypagestyle{plain}{
  \fancyhf{}
  \renewcommand{\headrulewidth}{0pt}
  \renewcommand{\footrulewidth}{0pt}
  \fancyfoot[LE,RO]{\thepage}
}

% 정리 환경
\newtheoremstyle{bookdefinition}
  {8pt}{8pt}{\normalfont}{}{\bfseries}{}{0.7em}
  {\thmname{#1}\thmnumber{ #2}\thmnote{ (#3)}}
\newtheoremstyle{bookplain}
  {8pt}{8pt}{\itshape}{}{\bfseries}{}{0.7em}
  {\thmname{#1}\thmnumber{ #2}\thmnote{ (#3)}}
\theoremstyle{bookdefinition}
\newtheorem{definition}{정의}[chapter]
\newtheorem{example}{예제}[chapter]
\theoremstyle{bookplain}
\newtheorem{theorem}{정리}[chapter]
\newtheorem{proposition}[theorem]{명제}
\newtheorem{lemma}[theorem]{보조정리}
\renewcommand{\proofname}{증명}

% 상자
\newtcolorbox{objectivebox}[1][학습 목표]{
  enhanced,breakable,
  colback=BookBlue!4,colframe=BookBlue,boxrule=0.6pt,arc=1mm,
  title={#1},fonttitle=\sffamily\bfseries,coltitle=white,colbacktitle=BookBlue,
  attach boxed title to top left={xshift=5mm,yshift*=-\tcboxedtitleheight/2},
  boxed title style={boxrule=0pt,arc=1mm},
  top=5mm,left=4mm,right=4mm,bottom=3mm
}

\tcbset{booksidebar/.style={
  enhanced,breakable,blanker,
  left=9mm,right=9mm,top=2mm,bottom=2mm
}}
\newtcolorbox{infobox}[1][정보]{
  booksidebar,
  borderline west={2.2pt}{0pt}{BookBlue},
  borderline east={2.2pt}{0pt}{BookBlue},
  before upper={\textcolor{BookBlue}{\bfseries 정보}\quad\textbf{#1}\par\smallskip}
}
\newtcolorbox{tipbox}[1][학습 팁]{
  booksidebar,
  borderline west={2.2pt}{0pt}{BookGreen},
  borderline east={2.2pt}{0pt}{BookGreen},
  before upper={\textcolor{BookGreen}{\bfseries 학습 팁}\quad\textbf{#1}\par\smallskip}
}
\newtcolorbox{cautionbox}[1][주의]{
  booksidebar,
  borderline west={2.2pt}{0pt}{BookOrange},
  borderline east={2.2pt}{0pt}{BookOrange},
  before upper={\textcolor{BookOrange}{\bfseries 주의}\quad\textbf{#1}\par\smallskip}
}
\newtcolorbox{warningbox}[1][경고]{
  booksidebar,
  borderline west={2.2pt}{0pt}{BookRed},
  borderline east={2.2pt}{0pt}{BookRed},
  before upper={\textcolor{BookRed}{\bfseries 경고}\quad\textbf{#1}\par\smallskip}
}

\newtcolorbox{shadequote}[1][]{
  enhanced,breakable,
  colback=black!4,colframe=black!18,boxrule=0pt,leftrule=2pt,arc=0pt,
  left=6mm,right=5mm,top=3mm,bottom=3mm,fontupper=\itshape,
  after upper={
    \if\relax\detokenize{#1}\relax\else
      \par\smallskip\hfill--- #1
    \fi
  }
}

\tcbset{bookexercise/.style={
  enhanced,breakable,boxrule=0pt,leftrule=3pt,arc=0pt,
  coltitle=white,fonttitle=\sffamily\bfseries,
  attach boxed title to top left={xshift=0pt,yshift*=-\tcboxedtitleheight},
  top=5mm,left=5mm,right=4mm,bottom=3mm
}}
\newtcolorbox{exercisebox}[1][연습문제]{
  bookexercise,colback=BookBlue!4,colframe=BookBlue,title={#1},
  boxed title style={colback=BookBlue,boxrule=0pt,arc=0pt}
}
\newtcolorbox{answerbox}[1][풀이]{
  bookexercise,colback=BookGreen!4,colframe=BookGreen,title={#1},
  boxed title style={colback=BookGreen,boxrule=0pt,arc=0pt}
}

% 코드
\lstdefinestyle{bookcode}{
  basicstyle=\ttfamily\small,
  backgroundcolor=\color{black!3},
  frame=single,
  rulecolor=\color{black!15},
  numbers=left,
  numberstyle=\scriptsize\color{BookGray},
  numbersep=6pt,
  xleftmargin=2.2em,
  framexleftmargin=1.7em,
  showstringspaces=false,
  breaklines=true,
  columns=fullflexible,
  keepspaces=true,
  aboveskip=8pt,
  belowskip=8pt
}
\lstnewenvironment{codeblock}[1][]{\lstset{style=bookcode,#1}}{}

% 참고문헌
\usepackage[backend=biber,style=authoryear,maxbibnames=99]{biblatex}
\addbibresource{references.bib}

% 링크
\usepackage{hyperref}
\hypersetup{
  unicode=true,colorlinks=true,
  linkcolor=black,citecolor=black,urlcolor=black,
  pdftitle={\BookTitle},pdfauthor={\BookAuthor}
}

% 제목 속 수식은 인쇄본에서는 그대로 두고 PDF 책갈피에서만 일반 문자열로 바꾼다.
\pdfstringdefDisableCommands{%
  \def\({}%
  \def\){}%
  \def\alpha{alpha}%
  \def\sigma{sigma}%
  \def\tau{tau}%
  \def\pi{pi}%
}

% 자주 쓰는 명령
\newcommand{\keyterm}[1]{\textbf{#1}\index{#1}}
\newcommand{\figref}[1]{그림~\ref{#1}}
\newcommand{\tabref}[1]{표~\ref{#1}}
\newcommand{\eqnref}[1]{식~\eqref{#1}}
\newcommand{\E}{\mathbb{E}}
\newcommand{\Pprob}{\mathbb{P}}
\newcommand{\Qprob}{\mathbb{Q}}
\newcommand{\R}{\mathbb{R}}
\newcommand{\ind}{\mathbf{1}}
\newcommand{\dd}{\mathop{}\!\mathrm{d}}
